\section{Schichtweise Entstehung mit dem Logit Lens}

Während das vorige Kapitel gezeigt hat, dass Sentiment im Embedding-Raum vorhanden ist, untersucht dieses Kapitel, ab welcher Schicht das Modell diese Information für seine Vorhersage nutzt. Dazu dient der Logit Lens.

\subsection{Verfahren}

Der Logit Lens beruht auf der Idee des Residual Streams \cite{nostalgebraist2020logit}. Da die finale Ausgabe des Modells eine Projektion des Residual Streams auf den Vokabularraum ist, lässt sich dieselbe Ausgabematrix auch auf frühere Hidden States anwenden. Für jede Schicht entstehen so Zwischen-Logits, die zeigen, welches nächste Token das Modell nach dieser Schicht bevorzugen würde. Ein wichtiger Vorbehalt besteht darin, dass die Ausgabematrix für die letzte Schicht trainiert wurde. Mittlere Schichten könnten Informationen anders kodieren, weshalb der Logit Lens einzelne Inhalte übersehen kann.

\subsection{Schichtweise Vorhersagen}

Für ein positives und ein negatives Prompt-Paar werden die wahrscheinlichsten Token je Schicht bestimmt. In den frühen Schichten von etwa 0 bis 8 dominieren unspezifische Token mit sehr geringen Wahrscheinlichkeiten. Ab den mittleren Schichten erscheinen erstmals Wörter aus dem Hu-\&-Liu-Lexikon. In den späten Schichten verstärken sich die sentimentbezogenen Vorhersagen deutlich. Tabelle~\ref{tab:logitlenstokens} fasst zentrale Beobachtungen zusammen.

\begin{table}[htbp]
  \centering
  \caption{Ausgewählte Beobachtungen des Logit Lens für das positive und das negative Prompt-Paar.}
  \label{tab:logitlenstokens}
  \begin{tabular}{lll}
    \toprule
    Beobachtung & Token & Wahrscheinlichkeit \\
    \midrule
    erstes positives Top-1-Token (Schicht 12) & lucky & sehr gering \\
    positives Top-1-Token (letzte Schicht) & satisfied & 0,1319 \\
    negatives Top-1-Token (letzte Schicht) & sick & 0,1604 \\
    \bottomrule
  \end{tabular}
\end{table}

Abbildung~\ref{fig:topkheat} zeigt die wahrscheinlichsten Token je Schicht für den positiven und den negativen Prompt. In den frühen und mittleren Schichten sind die Vorhersagen wenig interpretierbar und wechseln häufig. Ab etwa Schicht 17 steigen die Wahrscheinlichkeiten deutlich an, und die Vorhersagen konzentrieren sich auf wenige dominante Token. Abbildung~\ref{fig:huliumass} stellt dar, wie viel Wahrscheinlichkeitsmasse der Top-Token je Schicht auf Lexikonwörter entfällt.

\begin{figure}[htbp]
  \centering
  \includegraphics[width=0.49\textwidth]{logit_lens_positive_topk_heatmap.png}
  \hfill
  \includegraphics[width=0.49\textwidth]{logit_lens_negative_topk_heatmap.png}
  \caption{Wahrscheinlichste Token je Schicht für den positiven (links) und den negativen Prompt (rechts).}
  \label{fig:topkheat}
\end{figure}

\begin{figure}[htbp]
  \centering
  \includegraphics[width=0.49\textwidth]{logit_lens_positive_hu_liu_topk_mass.png}
  \hfill
  \includegraphics[width=0.49\textwidth]{logit_lens_negative_hu_liu_topk_mass.png}
  \caption{Anteil der Wahrscheinlichkeitsmasse auf Hu-\&-Liu-Wörtern je Schicht für den positiven (links) und den negativen Prompt (rechts).}
  \label{fig:huliumass}
\end{figure}

\subsection{Logit-Differenz als Sentiment-Score}

Um die Entwicklung über die Schichten zu quantifizieren, wird je Schicht ein Sentiment-Score gebildet. Dieser ergibt sich aus dem durchschnittlichen Logit-Wert über alle positiven Lexikonwörter abzüglich des durchschnittlichen Logit-Werts über alle negativen Lexikonwörter. Ein höherer Wert bedeutet eine stärkere Ausrichtung der Schicht auf positives Sentiment. Der Abstand zwischen dem Score des positiven und des negativen Prompts beschreibt, wie deutlich das Modell beide Eingaben trennt.

In den frühen Schichten liegen die Werte nahe null, positive und negative Eingaben besitzen also zunächst sehr ähnliche Repräsentationen. Auffällig ist dabei, dass die absoluten Sentiment-Scores beider Klassen zunächst nicht einfach konstant bleiben, sondern ab etwa Schicht 13 gemeinsam ins Negative absinken und bei beiden Prompt-Paaren zwischen Schicht 19 und 22 Werte von bis zu -2,8 erreichen. Sowohl positive als auch negative Lexikonwörter erhalten in diesem Bereich also einen unterdurchschnittlichen Logit-Wert; das Modell scheint die Logit-Masse zwischenzeitlich auf andere Token zu verteilen, bevor in der letzten Schicht ein Sprung auf positive Werte von bis zu +5,5 erfolgt. Der Abstand zwischen den beiden Klassen, also die eigentliche Sentiment-Trennung, wächst über diesen Bereich hinweg tendenziell an, jedoch nicht durchgehend monoton: Beim kurzen Prompt-Paar zeigt die Logit-Differenz bei Schicht 14 einen lokalen Spitzenwert, der bei Schicht 15 wieder zurückgeht, und schwankt bereits in den Schichten 0 bis 10 um die Null-Linie. Beim längeren Prompt-Paar verläuft die Differenz dagegen deutlich glatter und nahezu durchgehend steigend. Die größten Werte werden bei beiden Varianten erst zwischen Schicht 17 und Schicht 23 erreicht, der mit Abstand größte Einzelschritt liegt jedoch im Übergang von Schicht 23 zu Schicht 24, in dem sich der Abstand jeweils mehr als verdoppelt. Dies deutet darauf hin, dass die finale Schicht zusätzlich die konkrete Next-Token-Vorhersage optimiert. Kurzes und längeres Prompt-Paar zeigen damit einen ähnlichen Gesamttrend und denselben Schlusssprung, unterscheiden sich aber sichtbar in der Glattheit des Verlaufs zwischen den beiden Extremen. Die Abbildungen~\ref{fig:ldscores} und~\ref{fig:lddiff} zeigen die Verläufe der Sentiment-Scores und der Logit-Differenzen.

\begin{figure}[htbp]
  \centering
  \includegraphics[width=0.49\textwidth]{logit_lens_prompt_hu_liu_logit_scores_Trial1.png}
  \hfill
  \includegraphics[width=0.49\textwidth]{logit_lens_prompt_hu_liu_logit_scores_Trial2.png}
  \caption{Sentiment-Scores je Schicht für das kurze (links) und das längere (rechts) Prompt-Paar.}
  \label{fig:ldscores}
\end{figure}

\begin{figure}[htbp]
  \centering
  \includegraphics[width=0.49\textwidth]{logit_lens_prompt_hu_liu_logit_difference_Trial1.png}
  \hfill
  \includegraphics[width=0.49\textwidth]{logit_lens_prompt_hu_liu_logit_difference_Trial2.png}
  \caption{Logit-Differenz je Schicht für das kurze (links) und das längere (rechts) Prompt-Paar.}
  \label{fig:lddiff}
\end{figure}

\subsection{Aggregation über den CAD-Datensatz}

Damit das Ergebnis nicht von einzelnen Prompts abhängt, wird die Logit-Differenz über alle 1.707 gegenfaktischen Paare des CAD-Datensatzes aggregiert \cite{kaushik2020learning}. Abbildung~\ref{fig:cadaggregate} zeigt den gemittelten Verlauf zusammen mit der Streuung von einer Standardabweichung. In den frühen Schichten ist die mittlere Logit-Differenz nahezu null, und das Streuband reicht bis etwa Schicht 10 auf beiden Seiten an die Null-Linie heran. Für einen relevanten Teil der 1.707 Paare lässt sich Sentiment in diesen frühen Schichten demnach noch nicht zuverlässig trennen, auch wenn der Mittelwert bereits leicht positiv ist. Auch im Mittelwert zeigt sich bei Schicht 14 derselbe lokale Bump wie bei den beiden einzelnen Prompt-Paaren, gefolgt von einem Rückgang bei Schicht 15 -- das Muster ist also kein Artefakt eines einzelnen Beispiels, sondern systematisch über den Datensatz hinweg nachweisbar. Ab Schicht 16 steigt der Mittelwert dann weitgehend stetig an und erreicht sein Maximum mit dem Sprung von Schicht 23 auf Schicht 24, wobei auch die Streuung in den oberen Schichten deutlich zunimmt.

\begin{figure}[htbp]
  \centering
  \includegraphics[width=0.8\textwidth]{logit_lens_cad_logit_difference_aggregate.png}
  \caption{Über 1.707 CAD-Paare gemittelte Logit-Differenz je Schicht.}
  \label{fig:cadaggregate}
\end{figure}

Der Logit Lens beantwortet die zweite Forschungsfrage. Sentiment wird nicht in einer einzelnen Schicht gespeichert, sondern schrittweise über den Residual Stream aufgebaut. Die deutlichste Trennung zwischen positiven und negativen Eingaben entsteht in den oberen Schichten zwischen etwa Schicht 17 und Schicht 23. Dieser Befund stützt die Leithypothese und macht diesen Schichtbereich zum vorrangigen Kandidaten für die folgenden Analysen.
